\documentclass[11pt,a4paper]{article}
\usepackage[T1]{fontenc}
\usepackage[utf8]{inputenc}
\usepackage{newpxtext}
\usepackage[margin=25mm,headheight=15pt]{geometry}
\usepackage{amsmath,amsthm,mathtools}
\usepackage{newpxmath}
\usepackage{microtype,booktabs,array,longtable}
\usepackage[dvipsnames]{xcolor}
\usepackage{enumitem,fancyhdr,tcolorbox}
\usepackage[colorlinks=true,linkcolor=MidnightBlue,citecolor=MidnightBlue,
 urlcolor=MidnightBlue,pdftitle={Signed binomial supercongruences and OEIS A352373},
 pdfauthor={Research note prepared with ChatGPT}]{hyperref}
\usepackage[nameinlink,noabbrev]{cleveref}
\definecolor{Ink}{HTML}{21394D}
\definecolor{Accent}{HTML}{176B68}
\definecolor{Pale}{HTML}{F1F6F7}
\pagestyle{fancy}
\fancyhf{}
\fancyhead[L]{\small\color{Ink}Signed binomial supercongruences}
\fancyhead[R]{\small\color{Ink}Research note \textbullet\ 18 September 2026}
\fancyfoot[C]{\small\thepage}
\renewcommand{\headrulewidth}{0.3pt}
\setlength{\parindent}{0pt}
\setlength{\parskip}{4.5pt plus 1pt}
\linespread{1.035}
\setlength{\emergencystretch}{2em}
\numberwithin{equation}{section}
\newtheoremstyle{research}{7pt}{7pt}{\itshape}{}{
\color{Ink}\bfseries}{.}{0.5em}{}
\theoremstyle{research}
\newtheorem{theorem}{Theorem}[section]
\newtheorem{lemma}[theorem]{Lemma}
\newtheorem{proposition}[theorem]{Proposition}
\newtheorem{corollary}[theorem]{Corollary}
\theoremstyle{definition}
\newtheorem{remark}[theorem]{Remark}
\newcommand{\Z}{\mathbb Z}
\newcommand{\Q}{\mathbb Q}
\newcommand{\Zloc}{\mathbb Z_{(p)}}
\newcommand{\vp}{v_p}
\newcommand{\coeff}[2]{[x^{#1}]\,#2}
\newcommand{\CC}{\mathcal C}
\newcommand{\dd}{\delta_p}
\newcommand{\floor}[1]{\left\lfloor #1\right\rfloor}
\newcommand{\Li}{\operatorname{Li}}
\newcommand{\oeis}[1]{\href{https://oeis.org/#1}{\texttt{#1}}}
\newcommand{\status}{\textsc{Unrefereed research draft}}

\begin{document}
\thispagestyle{empty}
{\large\color{Accent}\textsc{Combinatorics \& arithmetic generating functions}\par}
\vspace{6mm}
{\huge\bfseries\color{Ink}Signed binomial supercongruences\par}
\vspace{2mm}
{\Large A proof of the conjecture recorded in OEIS A352373\par}
\vspace{5mm}
{\normalsize Research note prepared with ChatGPT\hfill 18 September 2026\par}
\vspace{4mm}
\begin{abstract}
The OEIS entry A352373 records a conjecture that every sequence
$a_{r,s}(n)=[x^n](1+x)^{rn}(1-x)^{sn}$, with arbitrary integer $r,s$,
satisfies the two-term congruences modulo $p^{3h}$ for primes $p\ge5$.
We give a proof, allowing the coefficient index to be any nonnegative
integer multiple $tn$. The argument separates indices divisible by $p$
from the remaining indices. Jacobsthal's binomial congruence treats the
first part; a multiscale alternating reciprocal-square cancellation treats
the second. The same proof gives modulus $3^{3h-1}$ at $p=3$.
We derive four OEIS specializations, a denominator restriction for a
M\"obius transform, and formal Lambert-series and trilogarithm expansions.
Exact computations check 13,690 family congruences and 15,154 proof-component
instances. The proof uses a precisely cited classical congruence; it has
not been independently refereed or formally verified, and publication
priority is not claimed.
\end{abstract}

\begin{tcolorbox}[colback=Pale,colframe=Accent,boxrule=0.6pt,arc=0pt,
 left=3mm,right=3mm,top=2mm,bottom=2mm]
\textbf{Main result.} For $r,s\in\Z$, $t\in\Z_{\ge0}$, and
$a(n)=[x^{tn}](1+x)^{rn}(1-x)^{sn}$,
\[
 \boxed{\quad a(mp^h)\equiv a(mp^{h-1})\pmod{p^{3h}}
 \quad(p\ge5).\quad}
\]
Here $m,h\ge1$ are arbitrary integers; in particular, $p\nmid m$ is
\emph{not} required. Negative exponents are interpreted as formal power
series. The conjecture in A352373 is exactly the case $t=1$.
\end{tcolorbox}

\section{The selected problem and the theorem}
\label{sec:problem}
The entry \oeis{A352373}, authored by Peter Bala, records both an individual
supercongruence conjecture and its two-parameter extension \cite{A352373}.
The individual sequence is
\begin{equation}\label{eq:selected}
 a(n)=[x^n]\frac{1}{(1-x)^{2n}(1-x^2)^n}
     =[x^n](1+x)^{-n}(1-x)^{-3n}.
\end{equation}
The general assertion concerns the same congruence for every integer pair
$(r,s)$ in $[x^n](1+x)^{rn}(1-x)^{sn}$. In the entry retrieved for this
note, that assertion is still explicitly marked as a conjecture.

Rather than finding a counterexample, the exact tests suggested a uniform
proof. We work with the larger family
\begin{equation}\label{eq:family}
 a_{r,s,t}(n):=[x^{tn}](1+x)^{rn}(1-x)^{sn},
 \qquad r,s\in\Z,\quad t,n\in\Z_{\ge0}.
\end{equation}
All of these coefficients are integers, including when $r$ or $s$ is
negative. Put
\begin{equation}\label{eq:delta}
 \delta_p=\begin{cases}1,&p=3,\\0,&p\ge5,\end{cases}
 \qquad\text{for odd primes }p.
\end{equation}

\begin{theorem}\label{thm:family}
For every odd prime $p$, all integers $r,s$, all nonnegative integers $t$,
and all positive integers $m,h$,
\begin{equation}\label{eq:main}
 a_{r,s,t}(mp^h)\equiv a_{r,s,t}(mp^{h-1})
                  \pmod{p^{3h-\delta_p}}.
\end{equation}
\end{theorem}

A homogeneous coefficient statement, proved in \cref{sec:proof}, implies
this theorem immediately. The proof uses the classical Jacobsthal
congruence in the version allowing negative integer upper entries
\cite[Lemma 5.1]{Straub}. Its organization follows the established
binomial-descent and block-cancellation approach to supercongruences
\cite[Section 2]{OSS}. The purpose here is the explicit application to the
whole coefficient family, not a claim that those proof techniques are new.

The distinction between a conjectural entry and a new theorem matters:
\oeis{A234839}, the case $(r,s,t)=(1,2,1)$, already has this
supercongruence, as recorded in the entry and in
\cite[Example 3.3]{OSS}. We do not count that special case as a new result
\cite{A234839}. The limited literature check underlying the present note
is described in \cref{sec:status} and in the accompanying source log.

\section{Elementary preliminaries and the classical input}
\label{sec:prelim}
For an integer $u$ and a nonnegative integer $k$, use
\[
 \binom uk=\frac{u(u-1)\cdots(u-k+1)}{k!},\qquad \binom u0=1.
\]
In particular, for $v\ge1$,
\begin{equation}\label{eq:negative}
 \binom{-v}{k}=(-1)^k\binom{v+k-1}{k}\in\Z.
\end{equation}
If $u\ge0$ and $k>u$, the coefficient is zero. We never need binomial
coefficients with negative lower entries in the proof.

Let $\Zloc$ be the ring of rational numbers whose reduced denominators
are not divisible by $p$. A congruence in this ring means divisibility by
the indicated power of $p$ inside $\Zloc$. Thus all reciprocal squares
$1/j^2$ with $p\nmid j$ are legitimate. We set $v_p(0)=+\infty$.
No convergence of infinite $p$-adic series will be needed.

\begin{lemma}[One-digit binomial descent]\label[lemma]{lem:descent}
Let $p$ be any prime, $h\ge1$, $p^h\mid U$, and $k\ge0$. Then
\begin{equation}\label{eq:descent}
 \binom{U-1}{k}\equiv
 (-1)^{k-\floor{k/p}}
 \binom{U/p-1}{\floor{k/p}}\pmod{p^h}.
\end{equation}
\end{lemma}
\begin{proof}
Splitting the product into indices divisible and not divisible by $p$
gives the exact identity in $\Q$
\[
 \binom{U-1}{k}
 =\binom{U/p-1}{\floor{k/p}}
   \prod_{\substack{1\le i\le k\\p\nmid i}}\frac{U-i}{i}.
\]
Each factor in the product is congruent to $-1$ modulo $p^h$ in $\Zloc$.
There are $k-\floor{k/p}$ such factors. The remaining binomial coefficient
is an integer, even for negative $U$, so multiplication preserves the
congruence. This argument also allows a vanishing remaining coefficient.
\end{proof}
This standard descent identity also appears as \cite[Lemma 2.4]{OSS}.

\begin{lemma}[Jacobsthal; imported classical theorem]\label[lemma]{lem:jacobsthal}
Let $p$ be an odd prime, let $X,Y$ be integers with $Y\ge0$ and $p\mid X,Y$, and
suppose $\binom{X/p}{Y/p}\ne0$. If $XY(X-Y)\ne0$, then
\begin{equation}\label{eq:jacobsthal}
 \frac{\binom XY}{\binom{X/p}{Y/p}}-1
 \in p^{v_p(XY(X-Y))-\delta_p}\Zloc.
\end{equation}
In the remaining nonvanishing cases $Y=0$ or $X=Y$, the ratio equals one.
\end{lemma}

This is the odd-prime part of \cite[Lemma 5.1]{Straub}: the exponent there
is $v_p(XY(X-Y)/12)$. Its loss is one at $p=3$ and zero at $p\ge5$.
That reference includes negative upper entries. The familiar
$p\ge5$ valuation form is also stated in \cite[Lemma 2.1]{OSS}.
We invoke this known theorem, rather than claim a new proof of it.

There is no use of a ratio $0/0$. If $\binom{X/p}{Y/p}=0$, then
$X\ge0$ and $Y>X$, and $\binom XY=0$ as well. Conversely, a zero fine
coefficient has a zero coarse coefficient. These support conditions are
unchanged by division of both entries by $p$.

\begin{lemma}[Alternating reciprocal-square blocks]\label[lemma]{lem:blocks}
Let $p$ be an odd prime, $\ell\ge1$, and $q\ge0$. Then
\begin{equation}\label{eq:blocks}
 H_{\ell,q}:=
 \sum_{\substack{qp^\ell\le j<(q+1)p^\ell\\p\nmid j}}
       \frac{(-1)^j}{j^2}
 \equiv0\pmod{p^\ell}.
\end{equation}
\end{lemma}
\begin{proof}
Write $P=p^\ell$. Within the block, pair $j=qP+a$ with
$j'=qP+P-a$, for $1\le a<P$ and $p\nmid a$. The pairing has no fixed
point since $P$ is odd. We have $j'\equiv-j\pmod P$ and
$(-1)^{j'}=-(-1)^j$. Their reciprocal-square contributions cancel modulo
$P$. All denominators involved are units modulo $P$.
\end{proof}

\section{The weighted cancellation that supplies the third power}
\label{sec:weighted}
The difficult part is not the unweighted congruence in
\cref{lem:blocks}, but retaining it in the presence of binomial weights.
The following lemma makes the necessary compatibility explicit. Its
telescoping step is a finite-block version of the weighted congruence
method; compare \cite[Lemma 5.6]{Straub}.

\begin{lemma}[Complementary-weight cancellation]\label[lemma]{lem:weighted}
Let $p$ be an odd prime, $h\ge1$, and $A,B,K\in\Z$ with $K>0$ and
$p^h\mid A,B,K$. Then
\begin{equation}\label{eq:weighted}
 W:=\sum_{\substack{1\le j<K\\p\nmid j}}
       \frac{(-1)^j}{j^2}
       \binom{A-1}{j-1}\binom{B-1}{K-j-1}
       \equiv0\pmod{p^h}.
\end{equation}
\end{lemma}
\begin{proof}
For $1\le\ell\le h$ and $0\le q<K/p^\ell$, define integer weights
\begin{equation}\label{eq:weights}
 U_\ell(q)=
 \binom{A/p^\ell-1}{q}
 \binom{B/p^\ell-1}{K/p^\ell-1-q}.
\end{equation}
All lower entries are nonnegative; negative upper entries cause no
integrality problem by \eqref{eq:negative}.

\textbf{The first descent has a minus sign.}
For $p\nmid j$ write $j=pq+a$ with $1\le a<p$. Because $p\mid K$,
\[
 \floor{(j-1)/p}=q,\qquad
 \floor{(K-j-1)/p}=K/p-1-q.
\]
Applying \cref{lem:descent} to both binomial coefficients, modulo $p^h$,
produces a sign with exponent
\[
 (j-1)+(K-j-1)-\bigl(q+(K/p-1-q)\bigr)=K-K/p-1.
\]
The difference $K-K/p=(p-1)K/p$ is even. Consequently,
\begin{equation}\label{eq:firstsign}
 \binom{A-1}{j-1}\binom{B-1}{K-j-1}
 \equiv- U_1(\floor{j/p})\pmod{p^h}.
\end{equation}
Introduce the sums
\begin{equation}\label{eq:T}
 T_\ell=
 \sum_{\substack{1\le j<K\\p\nmid j}}
 \frac{(-1)^j}{j^2}\,U_\ell(\floor{j/p^\ell}).
\end{equation}
Equation \eqref{eq:firstsign} says $W\equiv-T_1\pmod{p^h}$.

\textbf{Every subsequent descent has a plus sign.}
Suppose $1\le\ell<h$ and put $L=K/p^\ell$. Since $p\mid L$,
\[
 \floor{q/p}+\floor{(L-1-q)/p}=L/p-1.
\]
Both upper parameters $A/p^\ell$ and $B/p^\ell$ are divisible by
$p^{h-\ell}$. Another application of \cref{lem:descent} gives
\begin{equation}\label{eq:weightdescent}
 U_\ell(q)\equiv U_{\ell+1}(\floor{q/p})
                   \pmod{p^{h-\ell}}.
\end{equation}
Indeed, the combined sign exponent is $(L-1)-(L/p-1)=L-L/p$,
which is even. This is why the complementary lower entries in
\eqref{eq:weights} are useful.

\textbf{Block cancellation compensates for the declining precision.}
Group the difference $T_\ell-T_{\ell+1}$ by $q=\floor{j/p^\ell}$.
Since $K$ is a multiple of $p^\ell$, all blocks are complete, and
\begin{equation}\label{eq:telescoping}
 T_\ell-T_{\ell+1}
 =\sum_{q=0}^{K/p^\ell-1}
   \left(U_\ell(q)-U_{\ell+1}(\floor{q/p})\right)H_{\ell,q}.
\end{equation}
The bracket belongs to $p^{h-\ell}\Z$ by
\eqref{eq:weightdescent}, while $H_{\ell,q}\in p^\ell\Zloc$ by
\cref{lem:blocks}. Every product in \eqref{eq:telescoping} therefore
belongs to $p^h\Zloc$.

At the last level, each weight is constant on a block of length $p^h$:
\[
 T_h=\sum_{q=0}^{K/p^h-1}U_h(q)H_{h,q}\equiv0\pmod{p^h}.
\]
Thus $T_1\equiv T_2\equiv\cdots\equiv T_h\equiv0\pmod{p^h}$,
and $W\equiv0\pmod{p^h}$ follows. If $h=1$, only the first and last
steps are needed.
\end{proof}

\begin{remark}
The proof does not replace a weight by a constant modulo $p^h$ across an
entire block of length $p^h$ in a single step. At level $\ell$ it has only
$p^{h-\ell}$ precision; the block sum supplies exactly the missing
$p^\ell$. Ignoring this precision bookkeeping would leave a gap.
\end{remark}

\section{Proof of the coefficient supercongruence}
\label{sec:proof}
Define, for $A,B\in\Z$ and $K\ge0$,
\begin{equation}\label{eq:localcoefficient}
 \CC(A,B;K):=[x^K](1-x)^A(1+x)^B
 =\sum_{j=0}^K(-1)^j\binom Aj\binom B{K-j}.
\end{equation}
The finite convolution is valid for negative as well as positive
exponents in the ring of formal power series.

\begin{lemma}[Descent of the multiple indices]\label[lemma]{lem:multiples}
Let $p$ be an odd prime, $h\ge1$, $K\ge0$, and $p^h\mid A,B,K$.
For $0\le j\le K$ with $p\mid j$,
\begin{equation}\label{eq:multiples}
 \binom Aj\binom B{K-j}
 \equiv
 \binom{A/p}{j/p}\binom{B/p}{(K-j)/p}
 \pmod{p^{3h-\delta_p}}.
\end{equation}
\end{lemma}
\begin{proof}
If one coarse binomial coefficient is zero, its fine counterpart is
zero as well, so both products vanish. Assume otherwise. Put
\[
 e=\min\{h,v_p(j)\}.
\]
This is an integer between $1$ and $h$, with $e=h$ when $j=0$.
Both $j$ and $K-j$ are divisible by $p^e$. Applying
\cref{lem:jacobsthal} to each factor gives
\begin{equation}\label{eq:productratio}
 \binom Aj\binom B{K-j}
 =\binom{A/p}{j/p}\binom{B/p}{(K-j)/p}\,(1+E),
 \qquad E\in p^{h+2e-\delta_p}\Zloc.
\end{equation}
For example, $v_p(Aj(A-j))\ge h+2e$ whenever the three factors are
nonzero. Endpoint ratios equal one, so they satisfy the same error
bound without invoking a valuation of a nonzero product.

If $e<h$, then $v_p(j)=v_p(K-j)=e$, since $p^h\mid K$.
The identity
\[
 \binom uv=\frac uv\binom{u-1}{v-1}\qquad(v>0)
\]
shows that each coarse coefficient in \eqref{eq:productratio} is
divisible by $p^{h-e}$. Its upper entry is divisible by $p^{h-1}$,
and its lower entry has valuation $e-1$. If $e=h$, the same conclusion
means only that the coefficients are integers. Their product is
therefore always divisible by $p^{2(h-e)}$.

Subtracting the coarse product from \eqref{eq:productratio}, the
valuation of the difference is at least
\[
 (h+2e-\delta_p)+2(h-e)=3h-\delta_p.
\]
Since the difference is an integer, this proves the claimed ordinary
integer congruence.
\end{proof}

\begin{theorem}[Homogeneous coefficient theorem]\label{thm:homogeneous}
For an odd prime $p$, $h\ge1$, and integers $A,B,K$ with $K\ge0$ and
$p^h\mid A,B,K$,
\begin{equation}\label{eq:homogeneous}
 \CC(A,B;K)\equiv\CC(A/p,B/p;K/p)
                         \pmod{p^{3h-\delta_p}}.
\end{equation}
\end{theorem}
\begin{proof}
For $K=0$ both sides are one. Assume $K>0$ and split
\eqref{eq:localcoefficient} into the terms $p\mid j$ and $p\nmid j$.

For the first part, \cref{lem:multiples} applies term by term.
Because $p$ is odd, $(-1)^j=(-1)^{j/p}$ whenever $p\mid j$.
Consequently,
\begin{equation}\label{eq:regularsum}
 \sum_{\substack{0\le j\le K\\p\mid j}}
 (-1)^j\binom Aj\binom B{K-j}
 \equiv\CC(A/p,B/p;K/p)\pmod{p^{3h-\delta_p}}.
\end{equation}

For the second part, both $j$ and $K-j$ are units in $\Zloc$.
The exact binomial identity yields
\begin{equation}\label{eq:offexact}
 \binom Aj\binom B{K-j}
 =\frac{AB}{j(K-j)}\binom{A-1}{j-1}\binom{B-1}{K-j-1}.
\end{equation}
Also,
\[
 \frac1{j(K-j)}+\frac1{j^2}
   =\frac K{j^2(K-j)}\in p^h\Zloc.
\]
As $AB\in p^{2h}\Z$, summing \eqref{eq:offexact} over the unit indices
with signs gives
\begin{equation}\label{eq:offsum}
 \sum_{\substack{1\le j<K\\p\nmid j}}
 (-1)^j\binom Aj\binom B{K-j}
 \equiv -AB\,W\equiv0\pmod{p^{3h}},
\end{equation}
where the last congruence is \cref{lem:weighted}. If $A$ or $B$ is
zero, \eqref{eq:offexact} remains valid and the entire sum is zero.
Combining \eqref{eq:regularsum} and \eqref{eq:offsum} proves
\eqref{eq:homogeneous}.
\end{proof}

\begin{proof}[Proof of \cref{thm:family}]
Take $A=smp^h$, $B=rmp^h$, and $K=tmp^h$ in
\cref{thm:homogeneous}. By definition,
$\CC(A,B;K)=a_{r,s,t}(mp^h)$, and its scaled counterpart is
$a_{r,s,t}(mp^{h-1})$.
\end{proof}

\begin{remark}[Additional common divisibility]
If $m=p^v m_0$ with $p\nmid m_0$, the same theorem at height $h+v$
gives the stronger modulus $p^{3(h+v)-\delta_p}$. More generally,
\cref{thm:homogeneous} permits any common power of $p$ dividing all
three parameters $A,B,K$.
\end{remark}

A concrete check illustrates why cancellation is indispensable. For
\oeis{A352373} at $n=5$, the local parameters are
$(A,B,K)=(-15,-5,5)$. The $j=1$ summand is $1050=42\cdot5^2$,
not a multiple of $5^3$. The sum of the four nonmultiple-index terms is
$-8250=-66\cdot5^3$. Thus the extra power comes from their sum, not
from an unjustified termwise estimate.

\section{Consequences for the selected OEIS entries}
\label{sec:oeis}
The following substitutions are immediate from $1-x^2=(1-x)(1+x)$.
Each of these four entries explicitly records the $p^{3h}$
congruence as conjectural in the version consulted for this note
\cite{A352373,A348410,A351856,A351857}.
\begin{center}
\renewcommand{\arraystretch}{1.16}
\begin{tabular}{@{}lrrrc@{}}
\toprule
Sequence & $r$ & $s$ & $t$ & Defining coefficient \\
\midrule
\oeis{A352373} & $-1$ & $-3$ & $1$ & $[x^n](1-x)^{-2n}(1-x^2)^{-n}$ \\
\oeis{A348410} & $-1$ & $-2$ & $1$ & $[x^n](1-x)^{-n}(1-x^2)^{-n}$ \\
\oeis{A351856} & $-1$ & $-2$ & $2$ & $[x^{2n}](1-x)^{-n}(1-x^2)^{-n}$ \\
\oeis{A351857} & $-2$ & $-4$ & $1$ & $[x^n](1-x)^{-2n}(1-x^2)^{-2n}$ \\
\bottomrule
\end{tabular}
\end{center}
The coefficient-index extension $t\ge0$ is particularly useful for
\oeis{A351856}; its displayed coefficient representation requires $t=2$, beyond the
index used in the OEIS family conjecture.

For the selected sequence \eqref{eq:selected}, expanding the two
negative-binomial factors also gives the positive sum
\begin{equation}\label{eq:positivesum}
 a(n)=\sum_{k=0}^{\floor{n/2}}
       \binom{3n-2k-1}{n-2k}\binom{n+k-1}{k},\qquad n\ge1.
\end{equation}
For instance $a(1)=2$, $a(3)=74$, $a(5)=3252$, and $a(7)=154352$.
The theorem proves both
\[
 a(mp^h)\equiv a(mp^{h-1})\pmod{p^{3h}}\quad(p\ge5)
\]
and the additional relation
\[
 a(m3^h)\equiv a(m3^{h-1})\pmod{3^{3h-1}}.
\]
The latter assertion is a consequence of the weaker classical
Jacobsthal exponent at $p=3$, not an inference from finite examples.

\section{A generating-function consequence}
\label{sec:generating}
For any fixed $(r,s,t)$ let $a_n=a_{r,s,t}(n)$ and define
\begin{equation}\label{eq:mobius}
 b_n=\frac1{n^3}\sum_{d\mid n}\mu(d)a_{n/d},\qquad n\ge1,
\end{equation}
where $\mu$ is the M\"obius function.

\begin{proposition}\label[proposition]{prop:mobius}
For every $n\ge1$,
\begin{equation}\label{eq:denominator}
 b_n\in\Z[1/6],\qquad\text{and more precisely}\qquad
 3b_n\in\Z[1/2].
\end{equation}
In particular, $b_n$ is an integer when $\gcd(n,6)=1$.
The following identities hold as formal power series over $\Q$:
\begin{align}
 \sum_{n\ge1}a_nz^n
   &=\sum_{n\ge1}n^3b_n\frac{z^n}{1-z^n},\label{eq:lambert}\\
 \sum_{n\ge1}\frac{a_n}{n^3}z^n
   &=\sum_{n\ge1}b_n\Li_3(z^n),
 \qquad \Li_3(z)=\sum_{k\ge1}\frac{z^k}{k^3}.\label{eq:trilog}
\end{align}
\end{proposition}
\begin{proof}
Fix an odd prime divisor $p$ of $n$ and write $n=p^h m$ with
$p\nmid m$. Terms for which $p^2\mid d$ vanish in the M\"obius sum,
so its numerator equals
\[
 \sum_{e\mid m}\mu(e)
 \left(a_{(m/e)p^h}-a_{(m/e)p^{h-1}}\right).
\]
By \cref{thm:family}, this is divisible by $p^{3h-\delta_p}$.
Dividing by $n^3$ proves $v_p(b_n)\ge0$ for $p\ge5$, and
$v_3(b_n)\ge-1$. The reduced denominator of $b_n$ divides $n^3$,
so no other odd primes can occur. This proves \eqref{eq:denominator}
and the assertion for $\gcd(n,6)=1$.

M\"obius inversion of \eqref{eq:mobius} gives
$a_n=\sum_{d\mid n}d^3b_d$. The coefficient of $z^n$ on the right of
\eqref{eq:lambert} is precisely this divisor sum. The coefficient on
the right of \eqref{eq:trilog} is $n^{-3}$ times the same divisor sum.
For each degree only finitely many divisors contribute, so no analytic
convergence assumption is required.
\end{proof}

For \oeis{A352373}, the first seven $b_n$ are
\[
 2,\quad \frac54,\quad\frac83,\quad\frac{59}{8},\quad26,
 \quad\frac{308}{3},\quad450.
\]
This illustrates both the integrality away from $2,3$ and the genuine
possibility of a factor $3$ in the denominator. The transformation itself
is ordinary M\"obius inversion; the content supplied by the theorem is
its denominator restriction.

\section{Sharpness and failed generalizations}
\label{sec:boundaries}
The chosen conjecture survives; the following are counterexamples only
to \emph{stronger or broader} statements.

\textbf{The uniform exponent cannot be increased.}
For the selected sequence,
\[
 a(5)-a(1)=3252-2=3250=26\cdot5^3.
\]
It is not divisible by $5^4$. Thus a universal modulus $p^{3h+1}$,
and a fortiori $p^{4h}$, is false even within \oeis{A352373}.
This example establishes uniform sharpness; it does not assert that
equality of valuations holds for every prime or every height.

\textbf{The full exponent cannot be extended to $p=3$ or $p=2$.}
The same sequence gives
\[
 a(3)-a(1)=72=8\cdot3^2,\qquad a(2)-a(1)=10.
\]
The first is not divisible by $3^3$ and shows that the one-power loss
at $p=3$ is necessary in a uniform statement. The second is not
divisible by $2^3$. No general $2$-adic strengthening is claimed here.

\textbf{An arbitrary integer polynomial does not suffice.}
Consider instead $c(n)=[x^n](1+x+x^2)^n$. Direct expansion gives
$c(1)=1$ and
\[
 c(5)=\binom50\binom00+\binom52\binom21+
       \binom54\binom42=1+20+30=51.
\]
Consequently $c(5)-c(1)=50$ is not divisible by $5^3$. Even a
palindromic quadratic with constant coefficient one need not satisfy
the theorem. The opposite signs in $(1-x)^A(1+x)^B$ enable the specific
alternating cancellation used in the proof.

\section{Reproducible computation and proof audit}
\label{sec:computation}
The archive contains a standard-library-only Python program,
\texttt{code/verify.py}, together with its complete output
\texttt{results/verification.json} and a row-level record
\texttt{results/family\_cases.csv}. Run
\begin{verbatim}
python code/verify.py --output results
\end{verbatim}
The recorded run used Python 3.13.5. All checks passed.

\begin{center}
\begin{tabular}{@{}lr@{}}
\toprule
Test collection & Checks \\
\midrule
Parameter grid, primes at least five & 9,295 \\
Parameter grid, prime three with exponent $3h-1$ & 4,225 \\
Seeded additional family tests & 150 \\
Focused higher-height tests & 20 \\
\midrule
Total family congruences & 13,690 \\
Individual proof-component congruences & 15,154 \\
M\"obius-transform denominator checks & 500 \\
\bottomrule
\end{tabular}
\end{center}
The grid uses $-6\le r,s\le6$ and $0\le t\le4$, with the explicitly
listed $(p,h,m)$ cases stored in the JSON file. The focused tests include
height four at $p=5$, height five at $p=3$, and indices as large as
$11^3=1331$. There are 10,050 family cases with a nonzero difference;
the other 3,640 include legitimate zero or constant cases and are not
hidden in the totals. Randomized tests have the fixed seed 352373.

Two independent exact evaluators are compared whenever a sequence value
is computed. The first expands the binomial convolution. For the second,
write $F(x)=(1+x)^R(1-x)^S=\sum c_kx^k$. Logarithmic differentiation
and coefficient comparison give
\begin{equation}\label{eq:recurrence}
 (1-x^2)F'(x)=((R-S)-(R+S)x)F(x),\qquad
 k c_k=(R-S)c_{k-1}+(k-2-R-S)c_{k-2},
\end{equation}
with $c_{-1}=0$ and $c_0=1$. This uses $O(K)$ integer-arithmetic steps to
obtain $c_K$, with varying integer sizes. Each division is checked to be
exact. It is \emph{not} a modular recurrence that divides by a possible
nonunit $k$.

The proof-component tests include the first-descent minus sign, every
subsequent weighted transition, the terminal blocks, support-zero cases,
and the stronger $p^{3h}$ vanishing of the off-multiple sum even at $p=3$.
They test arithmetic instances of the proof, not just its conclusion.
They are still finite tests, not a formal verification of the universally
quantified theorem.

For a further transparent check, the selected sequence gives the
following valuations; full differences are recorded in the JSON file.
\begin{center}
\begin{tabular}{@{}rrrr@{}}
\toprule
$p$ & $h$ & $v_p(a(p^h)-a(p^{h-1}))$ & Guaranteed exponent \\
\midrule
5 & 1 & 3 & 3 \\
5 & 2 & 6 & 6 \\
5 & 3 & 9 & 9 \\
7 & 1 & 3 & 3 \\
3 & 1 & 2 & 2 \\
3 & 2 & 5 & 5 \\
\bottomrule
\end{tabular}
\end{center}

\section{Status, attribution, and what has actually been established}
\label{sec:status}
This note supplies an all-parameter argument for
\cref{thm:family}, not merely evidence for a small case. The only
non-elementary imported result is the classical congruence in
\cref{lem:jacobsthal}; all subsequent deductions, including the full
weighted cancellation, are given above. No unproved auxiliary
conjecture is used.

The source check covered the relevant OEIS entries and the two primary
papers cited for the congruence machinery. The general assertion remains
labelled conjectural in A352373 as retrieved on 18 September 2026.
An already-proved special case was explicitly separated from the target.
This is not an exhaustive literature review, and an equivalent general
result could exist in a different formulation. Accordingly, the claim is
\emph{a proof of the recorded assertion and its stated extensions}, not
verified priority for a new discovery.

The document is an unrefereed AI-assisted research draft. Neither an
independent mathematician's review nor a proof-assistant check was
performed. The arithmetic tests are reproducible but cannot replace
such review. The accompanying \texttt{PROOF\_AUDIT.md} identifies the
valuation, sign, boundary, and dependency checks that are especially
important for independent scrutiny. The proposed OEIS note in the
archive is a draft only; nothing has been submitted to OEIS.

\clearpage
\begin{thebibliography}{9}
\bibitem{A352373}
P.~Bala and OEIS contributors.
\emph{A352373: $a(n)=[x^n]((1-x)^{-2}(1-x^2)^{-1})^n$}.
The On-Line Encyclopedia of Integer Sequences.
\url{https://oeis.org/A352373}.
Entry author date: 14 March 2022. Consulted 18 September 2026.

\bibitem{A348410}
C.~E.~Lozada and OEIS contributors.
\emph{A348410: nonnegative integer solutions with an even-coordinate
restriction}.
The On-Line Encyclopedia of Integer Sequences.
\url{https://oeis.org/A348410}. Consulted 18 September 2026.

\bibitem{A351856}
P.~Bala and OEIS contributors.
\emph{A351856: the coefficient $[x^{2n}]((1-x)(1-x^2))^{-n}$}.
The On-Line Encyclopedia of Integer Sequences.
\url{https://oeis.org/A351856}. Consulted 18 September 2026.

\bibitem{A351857}
P.~Bala and OEIS contributors.
\emph{A351857: the coefficient $[x^n]((1-x)(1-x^2))^{-2n}$}.
The On-Line Encyclopedia of Integer Sequences.
\url{https://oeis.org/A351857}. Consulted 18 September 2026.

\bibitem{A234839}
V.~Kotesovec and OEIS contributors.
\emph{A234839: $\sum_{k=0}^n(-1)^k\binom nk\binom{2n}k$}.
The On-Line Encyclopedia of Integer Sequences.
\url{https://oeis.org/A234839}. Consulted 18 September 2026.

\bibitem{OSS}
R.~Osburn, B.~Sahu, and A.~Straub.
\emph{Supercongruences for sporadic sequences}.
Proceedings of the Edinburgh Mathematical Society (2) \textbf{59}
(2016), no.~2, 503--518.
\href{https://doi.org/10.1017/S0013091515000255}{doi:10.1017/S0013091515000255}.
\href{https://arxiv.org/abs/1312.2195v2}{arXiv:1312.2195v2} (2014).

\bibitem{Straub}
A.~Straub.
\emph{Multivariate Ap\'ery numbers and supercongruences of rational functions}.
Algebra \& Number Theory \textbf{8} (2014), 1985--2008.
\href{https://doi.org/10.2140/ant.2014.8.1985}{doi:10.2140/ant.2014.8.1985}.
\href{https://arxiv.org/abs/1401.0854v2}{arXiv:1401.0854v2}.
\end{thebibliography}
\end{document}
